\subsection{Wozu Computer?}
%\begin{wrapfigure}{l}{\linewidth}
 % \centering
 \begin{center}
 \includegraphics[width = 0.75\linewidth]{bilder/comics/otto1_1.png}
 \end{center}
%\end{wrapfigure}
 % \hspace{0cm}
  %\vspace{3cm}
\subsubsection{Vorlesungen Online}
Zu den meisten Vorlesungen kann man die Skripte im Internet finden. Das Institut 
für Informationssysteme (IFIS), mit dem ihr aber zu Beginn des Studiums 
weniger zu tun habt, stellt seine Vorlesungen sogar als Video zum Download bereit.

Hier die offiziellen Seiten zu den Vorlesungen:

% TODO das sind aber nicht alle, oder? Was ist mit dem Rest?
\begin{description}
\item[Algorithmen und Datenstrukturen]~\\
{\footnotesize\nurl{http://www.ibr.cs.tu-bs.de/courses/ws0910/aud/index.html}}
\item[Diskrete Mathematik]~\\
{\footnotesize\nurl{http://www.mathematik.tu-bs.de/dm/mitarb/kemnitz/diskrmathinf.html}}
\item[Programmieren I]~\\
{\footnotesize\nurl{http://www.ips.cs.tu-bs.de/content/view/260/198/lang,german/}}
\end{description}

Es gibt auch immer engagierte Studierende, die ihre Vorlesungsmitschriften online 
stellen. Da diese sehr wahrscheinlich in deinem Semester sind, hilft es, wenn du dich
in den Vorlesungen umhörst. Ansonsten ist \url{https://www.clevershit.de} die richtige
Anlaufstelle für den Informationsaustausch zwischen Studenten.

%\begin{figure}[h]
 \begin{center}
  \centering\includegraphics[width = 0.7\linewidth]{bilder/comics/otto1_2.png}
%  dddd
\end{center}
\subsubsection{Organisatorisches ohne Papier}
Ansonsten gibt es eine Reihe von Informationen, die ihr nur über das Web bekommt, 
und mehr und mehr Formalitäten (z.B. die Prüfungsanmeldung) werden auch in
die virtuelle Relatität verlagert.

Desweiteren k"onnt ihr euch im Netz euren individuellen Stundenplan
zusammenstellen, in Erfahrung bringen, wann die n"achsten Klausuren
stattfinden, lesen, was es in der Mensa zu essen gibt, endlich herausfinden,
wann das Fachbereichssekretariat ge"offnet hat, den Seminarraum des
Rechenzentrums f"ur Lerngruppen reservieren, offene HiWi"~Stellen bei den
Instituten finden und vieles mehr.

Die Webseiten der TU sind ein großer Dschungel, durch den man sich am besten 
mit Machete und Googlesuche kämpft. Um an der TU etwas zu finden, solltest 
du deinem eigentlichen Suchbegriff wahlweise "`tu braunschweig"' oder 
"`site:tu-braunschweig.de"' anhängen, und schon hast du gute Chancen zum Ziel 
zu kommen.

\begin{center}%h]
%  \centering
\includegraphics[width = 0.7\linewidth]{bilder/comics/otto1_3.png}
\end{center}
\subsubsection{Mitschreiben am PC}
Auf den ersten Blick mag es naheliegen, sich während der Vorlesungen Notizen 
am Laptop anzufertigen. In der Praxis gibt es da aber eine Reihe von Problemen, 
vor denen wir euch warnen möchten. Es hat schließlich seinen Grund, das nur 
10\% der Studenten in der Vorlesung am Laptop sitzen und davon 90\% diesen nur 
nutzen, um Farmville und Co. zu spielen: Die meisten Tafelanschriften bestehen 
aus verschachtelten Formeln, fremdartigen Buchstaben und verworrenen 
Zeichnungen. Diese in Echtzeit in den Laptop einzuhacken ist eine besondere
 Kunst, die du mit Notepad und Word gar nicht erst probieren brauchst. Eine 
 Chance hast du vielleicht mit einem Tablet PC, oder wenn du \LaTeX\ bereits 
 im Schlaf beherrscht - aber wer tut das schon zu Beginn des Studiums?

In den Vorlesungen, in denen du nicht Tafelweise abschreiben musst, sondern 
nur hier und da mal etwas notieren, zeigt sich der PC schon als nützlicher. 
Wenn du ab und zu den Vortrag des Profs damit vergleichen möchtest, was er 
in sein Skript geschrieben hat, kann dir der mitgebrachte Laptop unter 
Umständen das Ausdrucken von ein paar hundert Seiten ersparen. Du wirst aber 
schnell merken, dass es in praktisch keinem der Hörsääle und Seminarräume 
Steckdosen gibt, und in manchen nichtmals ausreichende WLAN-Signalstärke.

\subsubsection{Hausaufgaben am PC}
In vielen Fächern musst du regelmäßig Hausaufgaben erledigen und einreichen.
Keiner erwartet von dir, dass diese mit dem PC gemacht werden, es ist also
völlig ok sie von Hand zu schreiben. Es hat aber auch gewisse Vorteile, sie
am PC zu schreiben (z.B. mittels \LaTeX) und dann auszudrucken. Es gibt in
jedem Semester \LaTeX-Kurse die du im Schlüsselqualifiktationsbereich 
anrechnen lassen kannst, aber mit den Infos im WWW kann man sich das auch
selbst beibringen. Je eher du damit anfängst, desto weniger Probleme hast
du später, wenn du damit \footnote{Andere Programme gehen auch, aber
  wozu?} deine Abschlussarbeit aufsetzt.

\subsection{Computer-Pools an der Uni}
Es ist immer n"utzlich zu wissen, wo man mal schnell an einen Computer kann.
Zumindest ab und zu wirst du die Computer in der Uni benutzen, besonders
die Linuxarbeitspl"atze in \nroom{PK4.5} oder \nroom{PK4.8}, an denen du die Hausaufgaben f"ur
Programmieren abgeben musst.

{\small
An dieser Stelle gleich noch ein Tipp zu den Hausaufgaben in Programmieren:
Gib sie fr"uh ab, auch wenn du mehrere Wochen Zeit hast. In der letzten Woche
vor dem Abgabetermin ist der Raum immer so mit Studierenden und halbfertigen
Programmen "uberf"ullt, dass du oft sehr lange auf einen freien Computerplatz
warten musst - und noch l"anger auf den (dann schon total genervten) HiWi.
}

\begin{itemize}

\item[*] Im Erdgeschoss des Altbaus gibt es auf der rechten Seite zwei
Computerr"aume, einen weiter vorne (\nroom{PK4.6}) und einen genau in der Ecke
des Geb"audes (\nroom{PK4.5}). Zwei weitere R"aume (\nroom{PK4.8} und die
\nroom{"`Datenstation"'}) findest du im ersten Stock des Altbaus, auch wieder in
der rechten Ecke. Die Rechner in \nroom{PK4.5} und \nroom{PK4.8} sind mit Linux
ausgestattet. Im ersten Stock gibt es nun auch einen Windowsrechnerraum. Da kann man mal eine Word- oder
Powerpoint-Datei ausdrucken, wenn man denn muss.

\item[*] Reichlich Computer findet man schlie"slich im Gau"s-IT-Zentrum~(GITZ) an der
Hans-Sommer-Stra"se. Das ist der gedrungene, fast w"urfelf"ormige, dunkle
Klotz hinter dem Elektrotechnik-Hochhaus ("`E-Tower"'). Hier gibt es mehrere frei
zug"angliche R"aume mit Mappits\footnote{Das sind die ThinClients auf denen
Linux und Windows l"auft. Windows auf den Mappits befindet sich allerdings noch
im Teststadium.}, Linux PCs und auch einen gro"sen Windows"~Pool mit
Windows~2000 Rechnern. Es gibt hier auch Räume für Medienbearbeitung, wo du etwa
Video-Digitalisierer, ein Tonstudio und Rechner mit der Adobe Creative Suite 
Production Premium nutzen kannst.

\item[*] Seit 2010 stellt das IBR (Institut f"ur Betriebssysteme und Rechnerverbund)
im Raum G40 des Informatikzentrums einen Rechnerraum mit vielen, schnellen Linux-Rechnern 
zur Verf"ugung. Zu diesem CIP-Pool (Computer-Investitions-Programm) bekommt man mit 
seiner y-Nummer Zutritt. Wenn man Glück hat, funktioniert sogar einer der beiden Drucker
in diesem Raum, so dass man zum Drucken nicht das IZ verlassen muss.

\end{itemize}

% TODO gibt's die Mappits noch? Hab davon außerhalb der 1-te noch nie gehört. 
% Und ich habe im GITZ gelesen, man müsste seine MAC-Adresse freischalten lassen 
% um ins LAN zu kommen. Ist das noch aktuell?

\subsection{Der eigene Rechner}
Wenn du trotz aller Widrigkeiten planst, dir extra für's Studium einen 
(tragbaren) Rechner anzuschaffen, dann hast du hier gleich ein wenig 
Kaufberatung: Viel (Rechen- bzw. Grafik-)Leistung brauchst du im Studium 
nur für sehr wenige spezielle Fachgebiete - das einfachste Netbook wird also 
vermutlich schon reichen. Wichtiger ist vielleicht die Akkulaufzeit und die WLAN-Empfangsstärke.

\subsubsection{Welches System?}
Dir wird auffallen, dass zwar alle Systeme geduldet sind, aber die Linux hier
deutlich öfter über den Weg laufen wird als in der freien Wildbahn. Auch wir
sind große Linux-Fans und haben deshalb ab Seite \pageref{linux} ein paar Infos
dazu zusammengetragen.

Aber trotz dieser nicht ganz unauffälligen Beeinflussung gilt: Beim Betriebssystem 
hast du freie Wahl. Sämtliche Software, die du für's Studium brauchen 
könntest, gibt es für alle großen Systeme, meist sogar gratis. Für Linux ist eh 
praktisch alles frei erhältlich, für Windows spendiert 
Microsoft\footnote{Dafür zahlt die Uni wohl einen gewissen Betrag, aber für dich 
ist es gratis} den Studenten auch alles außer Office (siehe Seite \pageref{msdnaa}),
und auch Apple bringt dich dank satter Studentenrabatte durch Bachelor 
und Master. 

\subsubsection{Wege ins Uni-Netz}
Um den eigenen Rechner ins Netz zu bekommen, stehen die in der Uni WLAN und LAN 
offen. Zur Konfiguration siehe Seite \pageref{wlan}.

Für manche Aktivitäten (z.B. den Zugriff auf Prüfungsprotokolle) musst du dich 
direkt im Uni-Netz befinden. Wenn du und dein Rechner aber gerade zuhause oder sonstwo sind,
heißt dass nicht, dass du dich nun physisch auf den Weg machen musst. Mittels
VPN kannst du dich virtuell ins Uni-Netz einklinken. Schau einfach mal auf
den Seiten des GITZ nach, um mehr zu erfahren.

% Local Variables: 
% mode: latex
% TeX-master: "../../1-te"
% End: 
